\chapter{日辉族想要人类跪拜}

日辉族幸存者名叫索勒。

他是在争吵结束后自己说出来的。或许是因为所有人都看见了地球方案，继续只用族名互相称呼显得太空洞。又或许是法厄同正在碎裂，连骄傲也需要在死前找一个更具体的名字。

“索勒。”他说，“日辉第三律令塔下级执令官。”

瑟恩挑了挑眉：“终于从律令里长出名字了？”

索勒没有理他。

他站在地球投影前，残缺左臂垂在身侧，额骨上的太阳环裂出一道细缝。金色光从那缝里断续闪烁，让他看起来像一尊已经损坏却仍不肯倒下的神像。

“你们以为我们想要跪拜。”他说。

没有人回答。

因为他们确实这么以为。

索勒看着地球陆地上那些金色网格，声音低了些：“跪拜不是目的，是记忆方式。低阶生命无法理解星际尺度，也无法承受没有中心的恐惧。它们需要把高处具象化，需要天空、塔、王、神，需要知道有一个比自己更稳定的东西存在。”

“于是你们把自己放在那个位置。”瑟恩说。

“因为我们能承担。”

这句话让洛安心里生出一阵厌恶，却也有一点难以反驳的重量。

日辉族确实承担过很多东西。穹场最初由他们设计，法厄同北轨城市由他们维持，战争早期许多中立区靠他们的律令不至于立刻崩溃。对无数低阶族群来说，日辉族的秩序曾经意味着食物、航道、庇护和可预期的明天。

问题是，他们后来无法区分庇护和占有。

索勒抬手触碰投影。金色网格放大，未来地球上出现一座座阶梯城市。城市中心是高塔，塔顶连接天空观测器。人形生物跪在广场上，仰望光从塔顶落下。它们的身体还没有完全成为现代人类，却已经有了肩膀、双手、向上看的眼睛。

查理看得头皮发麻。

那些画面太熟悉了。

金字塔，神庙，王权，祭司，天命，太阳神，从天而来的教化者。

它们当然不会一模一样。这里呈现的是日辉族的设计草图，距离地球历史仍隔着漫长演化。可那种结构已经足够接近，接近到让人不舒服。

“你们想把未来人类养成仰望者。”洛安说。

索勒看向他：“仰望不是耻辱。”

“如果他们没有选择低头或不低头的权利呢？”

索勒沉默了一下。

“选择是成熟意识的奢侈。幼年文明首先需要方向。”

瑟恩冷笑：“然后等它成熟，你们会说它还不够成熟。”

索勒的脸色终于变了。

“你们震翼族永远只会说路。可路通向哪里？没有中心的道路只会让弱小者迷失。”

“至少迷失的人还有可能找到自己的方向。”

索勒刚想反驳，导航室再次剧烈震动。地球投影扭曲，日辉族金色方案忽然和其他族群方案重叠。阶梯城市下方出现鳞裔族地下库，广场边缘浮出节肢族并联节点，海洋里传来渊民族水声，天空中穿过震翼族路线。

所有方案被迫挤在同一颗地球上。

画面很混乱。

却也奇异地像未来。

查理忽然意识到，后世人类也许从未完全摆脱这些设计。每一族的碎片，都在他们身上留下了一点。人类会建塔，也会挖地库；会服从王权，也会反抗王权；会渴望集体，也会害怕失去自己；会从海里听见歌，也会一代代仰望天空。

索勒看着那混合画面，第一次露出迷茫。

“这样会很痛苦。”他说。

渊民族老人轻声回答：“也许生命本来就不是为了不痛苦。”

巨鲸种子低低震了一下。

地球投影中，那些跪拜的人影模糊了。他们站起，又跪下，又站起。动作不断重叠，像未来还没有决定自己要如何面对高处。

洛安看着那一幕，忽然觉得，比起跪拜或反抗，也许更重要的是，未来必须有机会自己决定。
